\section{Discussion}
\label{sec:discussion}

\subsection{What is now genuinely established}

After removing unnecessary overstatement, the paper supports four defensible claims.
First, it specifies a deterministic substrate ontology in which the microscopic degrees of freedom are the embedding variables of a 3D brane in $\R^4$.
Second, it cleanly separates the continuum long-wavelength description from the cubic lattice postulate that is meant to realize it microscopically.
Third, it identifies the geometric structures that necessarily appear if a narrowband polarization subspace is transported adiabatically.
Fourth, it formulates localized particle-like states as a concrete nonlinear bound-state problem rather than as an interpretive metaphor.

\subsection{Why the paper should stop short of stronger claims}

The cleaned-up version deliberately avoids claiming that the Standard Model, QCD, or general relativity have already been derived.
That stronger claim would require at least three things that are still missing:
a demonstrated localized spectrum, an extracted effective long-distance field theory with calibrated coefficients, and a clear account of multi-particle dynamics.
The benefit of the present restraint is that the paper becomes easier to evaluate.
Either the proposed substrate supports the required localized states and geometric diagnostics, or it does not.

\subsection{Why the proton/neutron program is the decisive next test}

The computational priority has shifted from an electron-first program to a baryon-first program.
This is not because the electron problem has become uninteresting, but because the current cubic-lattice hypothesis has an intrinsically three-channel internal sector.
A proton- or neutron-like search therefore tests the deepest new ingredient of the model---axis-aligned triplet mixing and its possible non-Abelian holonomy---instead of postponing it.
If the triplet sector cannot generate robust fm-scale bound states, then the interpretation of that sector as the seed of color-like dynamics loses most of its credibility.

\subsection{Open problems}

Several hard problems remain.
One must determine which constitutive laws actually support self-guided localization; quantify how much lattice anisotropy survives into the accessible regime; extract the effective coefficients of the slow envelope theory from the microscopic model rather than inserting them phenomenologically; and understand how long-range holonomy labels are related to interaction phenomenology.
These are substantial tasks.
But they are now separated more cleanly from what the manuscript already provides, which is exactly what a publishable theory-fragment paper should achieve.
